\selectlanguage{english}
Chapter XIV is a short experiment in reducing homotopy theory to a calculus of squares in a single category. Its eight manuscript pages introduce five proposed axioms and begin the examples. The editors' postface, which has no manuscript pagination, explains why the intended model-category example fails. The fragment is especially instructive because Grothendieck identifies the necessary finer setting before attempting the reduction: square diagrams must be localized as diagrams, together with their maps, rather than reconstructed from four arrows in the localized category.

The mathematical object under discussion can be described without ambiguity. For $\mathcal H=\mathcal MW^{-1}$, let $\operatorname{Cart}_{\mathrm{im}}$ consist of those commutative squares in $\mathcal H$ that are isomorphic to the underlying images of homotopy-cartesian strict squares in $\mathcal M$. With $\square$ denoting the walking commutative square, the relevant functor is
\[
\operatorname{dia}_{\square}:
\mathcal M^{\square}(W^{\square})^{-1}
\longrightarrow\mathcal H^{\square}.
\]
The domain retains coherent diagrams. The target retains ordinary commutativity after localization. Membership in $\operatorname{Cart}_{\mathrm{im}}$ is an existential statement about a coherent realization. Proposition~1 would promote it to a statement about every strict realization, and this is the false inference.

There is a specific historical reason for wanting the larger diagram theory. In \emph{Pursuing Stacks}, \S69, Grothendieck connects it with the trace and determinant problems that motivated him in the period of SGA~5; the letter to Thomason of 2~April~1991 again stresses derived diagram categories and mapping cones \citep{Grothendieck1983,Grothendieck1991Letter}. The present chapter tests whether one can retain a useful part of their calculus in a simpler language. The editors' preliminary notice suggests that recognition of the failed example explains the abrupt ending. That suggestion is conjectural; the postface establishes the mathematical failure, not the author's reason for interrupting the manuscript.

The formal axioms are consistent: any category with pullbacks, equipped with its ordinary pullback squares, satisfies them. Cart~(1) gives symmetry, Cart~(2) tests trivial squares against isomorphisms, and Cart~(3) combines pasting with cancellation. Cart~(4) requires a completion of every cospan. Cart~(5)(a) extends a map between the three-vertex boundaries of two distinguished squares; (b) asserts that every such extension is invertible whenever the boundary map is invertible. The commentary supplies the cube proof, including the stronger three-face hypothesis in the footnote, and proves that (b) already follows from Cart~(1)--(3).

The source's second corollary isolates the weak universal property that does survive in homotopy theory. For a homotopy pullback $P$ of $A\to C\leftarrow B$, there is a surjection
\[
[T,P]\longrightarrow[T,A]\times_{[T,C]}[T,B].
\]
A compatible pair in the target admits a homotopy between its composites; a lift remembers a choice of that homotopy. Passing to components forgets the choice, so the map need not be injective. In the loop square for $S^1$, with $T=S^0$ in pointed spaces, it is the map $\mathbb Z\to\{*\}$. The surjection is equivalent to Cart~(5)(a) for a class containing the constant squares. The much stronger assertion in (b) does not follow from this weak property alone.

Positive composition does hold for the homotopy-image class of every Quillen model category. The commentary proves this by applying fullness and conservativity of the arrow comparison twice, first to the model category and then to its arrow category, so that two cartesian witnesses acquire a single shared middle column. This is compatible with the postface: its counterexample concerns cancellation.

The editors' example distinguishes the underlying square from its realization. Factor the basepoint of a fibrant object as
\[
0\xrightarrow{\sim}C\twoheadrightarrow X,
\qquad P=C\times_X0.
\]
The resulting strict pullback is homotopy cartesian. Its image in $\mathcal H^{\square}$ is isomorphic to the square with apex $P$, base $X$, the other vertices zero, and all edges zero. But that zero-edge square, regarded strictly in $\mathcal M$, has comparison $0_P:P\to P$ to the fibrational replacement. It is homotopy cartesian exactly when $P$ is weakly zero. This proves the failure of Proposition~1 whenever the loop object is nontrivial. The same underlying square gives the two final counterexamples: a horizontal rectangle with left and outer squares distinguished but right square not distinguished, and a cube whose first three components are identities while the fourth is zero.

The worked algebraic example takes $X=k[0]$ and $C=(k\xrightarrow1k)$ in degrees $0,1$. Then $P=k[-1]$. A specified filling of the fixed zero-edge square is a scalar $\lambda$; the homotopy-pullback comparison is multiplication by $\lambda$ on $P$. Its cone is $k\xrightarrow\lambda k$, so the coherent square is cartesian precisely for $\lambda\ne0$. The commentary computes the transport through a contraction and its sign explicitly. It also checks the editors' topological qualification: for a CW complex, weak triviality of the loop object means contractibility of the base component, with no restriction on other components. For unbounded complexes in a Grothendieck abelian category it means acyclicity.

The literal axioms have a stronger limitation than the displayed counterexample alone reveals. Our development proves that, if $\mathcal H$ has binary products, the axioms force all ordinary pullbacks to exist and force the distinguished class to be exactly their class. The proof first compares a distinguished weak pullback with an existing ordinary one: a section and the invertibility of a cone-preserving idempotent make the comparison an isomorphism. This is the minimality argument familiar for weak kernels \citep[\S1, p.~6 of the 2014 manuscript]{Krause2016FlatCovers}. Pasting a distinguished square with a weak kernel pair then shows that pullbacks along monomorphisms are ordinary. Finally one pulls back the diagonal $C\to C\times C$ along $A\times B\to C\times C$. The complete argument is given in the marked development block; its application to these axioms was not located in the bounded literature search, and no priority is claimed.

Additive categories satisfy the binary-products hypothesis. Moreover, even without products, the literal loop object in a pointed category must be zero, because $0\to X\leftarrow0$ already has ordinary pullback zero. If the dual structure also exists, its suspension is zero. If the two classes coincide, the square with base $X$ and three zero vertices makes $X$ itself zero. Thus Grothendieck's additive coincidence speculation is false for the literal structures: ordinary pullbacks and pushouts in vector spaces give an immediate counterexample.

The modern positive statement belongs to coherent squares. For the inclusion $j:\Phi\hookrightarrow\square$ of the cospan, cartesian coherent squares are the essential image of the fully faithful right Kan extension $j_*$. Restriction consequently identifies their category with $\mathbb D(\Phi)$. This makes the completion functorial and proves the correct fourth-vertex assertion for coherent morphisms. Coherent rectangles admit pasting and cancellation. In a pointed derivator the suspension--loop adjunction exists; stability, rather than additivity, is the condition under which cartesian and cocartesian coherent squares agree, equivalently that adjunction is an equivalence \citetext{\citealp[Proposition~3.17]{Groth2013}; \citealp[Theorem~7.1, arXiv v2]{GrothPontoShulman2014}}. External derived Hom has the expected continuity, with the usual enrichment hypotheses for an internal version \citep[Proposition~6.13]{Cisinski2003}. None of these assertions makes the forgetful functor reflect the property disproved in the postface.

The proposed main example and Proposition~1 are therefore settled negatively, while the coherent replacements are established theorems. The manuscript's broad request for further compatibility axioms does not by itself specify a present open problem. This commentary does not classify all literal quadrangulated categories without binary products or produce a nonordinary example in that generality. Its audit keeps that limit separate from the explicit counterexamples and the proved structural deductions.

\medskip
\noindent\emph{Commentary, mathematical audit, and the proofs marked \textup{[ours]} were AI-produced by \mbox{GPT-6 Astra} (OpenAI, ChatGPT Pro mode), checked and integrated by Claude (Anthropic). The source is Grothendieck's Chapter XIV in the electronic edition deciphered and transcribed by M.~K\"unzer and edited by M.~K\"unzer, J.~Malgoire, and G.~Maltsiniotis, version 7~October~2010 \citep{Grothendieck1990XIV}.}\par
